% Generated draft: requires mathematical and source review.
\documentclass[11pt]{article}
\usepackage[margin=25mm]{geometry}
\usepackage{fontspec}
\setmainfont{Latin Modern Roman}
\usepackage{amsmath,amssymb,mathtools,unicode-math}
\setmathfont{Latin Modern Math}
\usepackage{enumitem,needspace,xurl,hyperref}
\setlength{\parindent}{0pt}
\setlength{\parskip}{5pt}
\setlength{\emergencystretch}{3em}
\newcommand{\N}{\mathbb{N}}
\newcommand{\Z}{\mathbb{Z}}
\newcommand{\Q}{\mathbb{Q}}
\newcommand{\R}{\mathbb{R}}
\newcommand{\C}{\mathbb{C}}
\newcommand{\F}{\mathbb{F}}
\newcommand{\A}{\mathbb{A}}
\newcommand{\PP}{\mathbb{P}}
\newcommand{\E}{\mathbb{E}}
\newcommand{\eps}{\varepsilon}
\newcommand{\dd}{\,\mathrm{d}}
\begin{document}
\section*{UnsolvedMath ID: 30005669}
\textbf{Source problem code:} catalogue entry
\subsection{Definitions and mathematical statement}
Let \(X\) be a noncompact Riemann surface, and write \(K_X\) for its holomorphic cotangent bundle. The conventional interpretation of the tuple in the source is
\[
\vec q=(q_2,\ldots,q_n),\qquad q_i\in H^0(X,K_X^{\otimes i}),
\]
with an integer \(n\geq2\). The excerpt does not explicitly specify this range or the degrees of the differentials, so this interpretation requires confirmation.

Write
\[
(E,\theta)=(K_{X,n},\theta(\vec q))
\]
for the source-designated associated Higgs bundle, and retain \(C_{K,X,n}\) for its designated compatibility structure. The supplied excerpt does not define these objects or the meaning of compatibility with \(C_{K,X,n}\); these are unresolved parts of the formulation, not additional freely chosen data.

For a holomorphic vector bundle \(E\) with Higgs field \(\theta\in H^0(X,\operatorname{End}(E)\otimes K_X)\), a smooth positive-definite Hermitian metric \(h\) is called harmonic if its Chern connection \(D_h\) satisfies
\[
F_{D_h}+[\theta,\theta^{*h}]=0.
\]
Here \(\theta^{*h}\) is the Hermitian adjoint Higgs field and the bracket is the graded commutator of endomorphism-valued forms. Define, provisionally,
\[
\mathcal H_C(X,\vec q)=\left\{h:\ h\text{ is a harmonic metric on }(K_{X,n},\theta(\vec q)),\quad h\text{ is compatible with }C_{K,X,n}\right\}.
\]
The compatibility condition in this display remains a source-dependent predicate whose definition must be recovered before the problem is fully specified.

\textbf{Questions.} For every noncompact Riemann surface \(X\) and every differential tuple \(\vec q\) admitted by the source's construction:
\begin{enumerate}
\item Is \(\mathcal H_C(X,\vec q)\neq\varnothing\)?
\item Can one specify a canonical notion of a best member of \(\mathcal H_C(X,\vec q)\), such that exactly one best member exists for every such \((X,\vec q)\)?
\end{enumerate}
The second question asks for a selection criterion as well as existence and uniqueness relative to that criterion. No criterion, normalization, boundary behavior, or equivalence relation for uniqueness is supplied. It therefore does not assert that all compatible harmonic metrics are equal.

\textbf{Qualification from the appended background.} The extracted main-results passage additionally assumes that \(X\) is hyperbolic, described there as equivalent to \(X\) being biholomorphic to neither \(\C\) nor \(\C^*\). It denotes by \(g_X\) the unique complete hyperbolic K\"ahler metric, with curvature normalization unspecified. It describes a reference metric \(h_X\) assembled from powers of \(g_X\) and fixed constants \(a_k\), chosen to make \(h_X\) harmonic for \((K_{X,n},\theta(0))\), followed by an increasing filtration \(0\subset F_1\subset\cdots\subset F_n\). The powers, bundle construction, and continuation of the filtration passage are damaged or truncated in the supplied extraction. They do not supply a reliable definition of the best metric. The hyperbolicity assumption belongs to that appended passage and the supplied existence assessment; the original questions themselves are stated for noncompact Riemann surfaces generally.
\subsection{Short English statement}
The problem asks whether every admitted tuple of holomorphic differentials on a noncompact Riemann surface gives an associated Higgs bundle with a compatible harmonic metric, and whether a canonical criterion can select exactly one such metric. The bundle construction, compatibility condition, and meaning of best are not defined in the supplied excerpt. Supplied dated curation reports existence on noncompact hyperbolic surfaces and uniqueness under boundedness or completeness conditions, while describing unrestricted uniqueness as unsettled; these claims require independent review.
\subsection{Sources}
\begin{itemize}
\item Supplied local UnsolvedMath registry record, permanent ID \(30005669\), source problem code \texttt{catalogue entry}, title \emph{Canonical Harmonic Metrics on Noncompact Riemann Surfaces}. The record lists proposed year 2023 and supplies no proposer attribution.
\item Supplied report citation: \emph{Teichm\"uller Theory: Classical, Higher, Super and Quantum} (2024), Oberwolfach Reports, pp.~1843--1892, DOI \texttt{10.4171/owr/2023/33}, \url{https://doi.org/10.4171/owr/2023/33}. The extracted original passage includes the page marker 1871. The report is supplied as the authoritative source; it has not been independently consulted for this draft.
\item Li--Mochizuki, \emph{Higgs bundles in the Hitchin section}, \url{https://arxiv.org/abs/2307.03365}. Supplied literature-review reference.
\item Li--Mochizuki, \emph{Complete solutions of Toda equations}, \url{https://arxiv.org/abs/2010.05401}. Supplied literature-review reference.
\item Literature assessment dated exactly \texttt{2026-08-22}: supplied dated curation, not independently verified. It labels the problem partially solved and reports existence of compatible harmonic metrics for arbitrary holomorphic differential tuples on noncompact hyperbolic surfaces, uniqueness under boundedness or completeness conditions, and unrestricted uniqueness as unsettled. The precise hypotheses and scopes of these reported results are not included.
\item The registry labels the statement \texttt{corrected\_verified} and records that the official report was checked directly. This is supplied provenance about an earlier verification, not a claim of source verification or mathematical review performed for this draft. The background itself requests independent verification of time-sensitive status and cited claims.
\end{itemize}
\end{document}
